%!TEX root = thesis.tex
\chapter{Design and Analysis}
\label{chap:design}

This chapter discusses the experimental framework used in this research and how it will evaluate database technologies for large-scale time series data management. The previous chapter considered the theoretical foundations of time series databases, their architectures and characteristics. This chapter builds on that to translate those insights into structured and reproducible experimental design. The goal is to define the way the comparative evaluation will be done in a way that is methodically sound, transparent and aligned with the requirements of the specific context of the research. 

The evaluation will focus on assessing the scalability and performance of the selected database systems under representative time series workloads. Instead of basing the research on abstract comparisons or isolated performance claims, this study adopts an experimental approach that is controlled and in which a number of database technologies are subjected to identical workloads under similar conditions. This ensures that any differences observed in the performance can be attributed to the architectural and implementation characteristics instead of inconsistencies in the setup of the experiment. 

To meet this goal, the chapter defines key components of the experimental design. It establishes its evaluation framework and methodological approach. This includes defining the scope of comparison, mentioning the relevant performance dimensions and outlining the principles that are used to guarantee fairness and reproducibility. Then, it describes how different database systems were selected to be included in the study alongside their rationale. Thirdly, it details the dataset design and workload models used to simulate realistic data generated by sensor-based systems. 

This chapter also specifies the benchmarking approach and the experimental environment include the configurations of the system, deployment conditions and hardware setup. These elements are necessary for making sure that performance measurements remain consistent and comparable across systems. It defines the metrics that are used to assess system performance, including query latency, ingestion throughput, storage efficiency, scalability behavior and resource utilization. Each metric is operationalized based on how it is measured and interpreted inside the experimental framework. 


Finally, the chapter outlines the experimental procedure and discusses the potential ways the design is limited. This includes considerations related to system configuration constraints, dataset realism and generalizability of the results. These factors are addressed to make sure the research is balanced and transparent. 


\section{Experimental Design Framework}
To evaluate database systems for time series data management, one needs a structured methodological approach that ensures comparability, consistency and reproducibility. To achieve this, this study adopts the evaluation framework proposed by \citet{Brown_1996}, which offers a systematic process for assessing competing technologies on the basis of measurable performance metrics. This framework is especially fitting for comparative system evaluation because of its emphasis on the identification of relevant performance dimensions, designing representative experiments and analyzing the observed differences between technologies. The framework is organized into three primary phases: descriptive modeling, experiment design and experiment evaluation. Every phase plays a critical and separate role in structuring the research process and making sure the evaluation stays grounded both in its theoretical understanding and empirical evidence. In the context of this study, the phases are mapped as follows:

The descriptive modeling phase sets up the relationship between the problem domain and the technologies being evaluated. It involves identifying the major characteristics of the workload and the system requirements which influence performance outcomes. As addressed in Chapter 2, the descriptive modeling showed the properties of time series data, the architectural features of database systems and the problems associated with sensor workloads. The insights from that chapter inform the selection of performance metrics and the design of the experimental workloads used in this research. 


The experiment design phase is the focus of this current chapter. It translates the conceptual understanding that was developed in the descriptive modeling into a concrete plan for experimentation. This phase will define the scope for evaluation, selected database systems and design the representative workloads and datasets. It will also specify the metrics to be used to assess performance. A key objective here is to make sure all systems get evaluated under consistent conditions to enable meaningful comparison. To achieve comparability, each database is subjected to a similar dataset, similar workload patterns and a similar hardware environment. That way, differences in performance can be attributed to system behavior instead of external factors. To do this, system configuration is considered carefully, including the indexing strategies, data ingestion methods and query execution settings. Where possible, default configurations mirror typical deployment scenarios while making sure no system is unfairly disadvantaged because of misconfiguration. 


During the experimental design phase, the formulation of evaluation criteria is considered. Instead of using just one performance metric, this study adopts a multi-dimensional evaluation approach which considers different aspects of system performance. These include query latency, which measures how responsive the system is to analytical queries, ingestion throughput, which is about the system’s ability to handle continuous data input and storage efficiency, which is about how effective data compression and storage utilization are. The research may also consider scalability, which is about the way performance changes with increases in data volume and resource utilization which offers insight into system efficiency under load. 


Throughout the experiment design, principles of transparency and repeatability are considered. All aspects of the setup including system configurations, hardware specifications and workload definitions are explicitly documented. This makes sure that the results can be reproduced under similar conditions, which is a fundamental requirement for empirical research. Besides, experimented are conducted a number of times to account for differences in system performance and the results are averaged to offer a more reliable representation of system behavior. 


The final phase of the \citet{Brown_1996} framework, experiment evaluation, involves executing the defined experiments, collecting the data and analyzing the results. This phase will be addressed in subsequent chapters where selected database systems are evaluated based on metrics defined in this chapter. The analysis is intended to identify performance trade-offs between systems and understand how differences in architecture influence observed behavior under time series workloads. The framework adopted strengthens the validity of the findings and offers a clear basis for interpreting the results in relation to the research questions defined earlier. 


\section{Selection of Database Systems}
For this study, the selection of database technologies is guided by the need to evaluate systems which represent distinct architectural approaches to time series data management while still staying practically deployable within the experimental environment scope. Instead of attempting to evaluate all possible time series databases, this study stays focused by selecting a few representative systems that allow for meaningful comparisons across different philosophies of design. Three systems are picked for the initial phase of the evaluation – MongoDB, TimescaleDB and InfluxDB. These systems were picked because together, they represent different paradigms for managing time series data – native time series database architectures, document-oriented storage and relational time-series extensions. This diversity ensures the study can examine how fundamentally different design decisions can influence system behavior under identical workload conditions. 

TimescaleDB is picked to represent relational database systems extended for time series workloads. It is built on top of PostgresSQL so it retains its relational model and SQL interface while adding time-based partitioning and optimizations for performance that are tailored to temporal data. It is included to allow the study to find out how far a relational system can adapt to handle time-stamped data without abandoning established principles. 

InfluxDB is added to the list as a purpose-built time series database intended to specifically serve high throughput ingestion and improve efficiency of time-based querying. It has a storage engine, built-in retention mechanism and query model which are optimized for temporal workloads. Investigating InfluxDB provides insight into the limitations and performance advantages of specialized systems which prioritize time series operations over general-purpose use. 

MongoDB represents document-oriented NoSQL approaches to time series data management. It has a flexible schema design and very recently introduced time series collections to allow it to accommodate heterogenous sensor data with diverse structures. It has been included in the study to provide insight into how a schema-flexible, general-purpose system performs when applied to structured time series workloads.


The selection of the above systems was also influenced by the practical considerations related to integration and implementation. All three databases offer mature Python libraries including pymongo for MongoDB, psycopg2 for TimescaleDB and influxdb-client for InfluxDB. These libraries make it possible to keep the interaction patterns consistent across systems within the experimental framework. Because there are stable client libraries, performance measurements can be conducted programmatically and reproducibly.

Notably, this study is designed to be extensible. While the initial evaluation will focus on these three systems, the framework for experimentation allows for the inclusion of more databases at a later stage. This modularity is made possible by standardized dataset formats, standardized measurement procedures and standardized definitions for the workloads. Consequently, future systems can be added into the evaluation without requiring major changes to the overall design. The goal of this selection is not to find one superior database system, but to examine how different architectural approaches work under similar workload conditions. A focus on a small but representative set of systems makes sure the evaluation remains both analytically meaningful and manageable. 


\section{Dataset Design and Characteristics}
How the dataset is designed in this study is critical to the experimental framework because it directly influences the relevance and validity of the study. Instead of relying on pre-existing benchmark datasets, this study creates a synthetic dataset that mirrors the characteristics of time series data from sensor deployments. That way, the study has precise control over data properties like structure, distribution and volume, making sure that the dataset aligns with the operational context defined earlier. 

The dataset is made to simulate data generated by a sensor network similar to the Smart Ocean platform. Every record in the set represents a time-stamped record produced by a sensor devise. Each record is structured to include a sensor identifier, timestamp field and many measurement attributes that correspond to operational or environmental variables. These attributes could include values like salinity, temperature, pressure or other sensor-specific metrics. This is done to reflect the multidimensional nature or data from real-world sensors. 


To support comparative evaluation across the different database systems, the dataset is defined in a way that can be mapped consistently to each model. For relational systems like TimescaleDB, the dataset is saved as a structured table with columns that are clearly defined. For InfluxDB, the same data is mapped to tags, measurements and fields with the different sensor identifiers represented as tags to support efficiency in filtering. In MongoDB, the records get stored like documents with field structures that are flexible to allow measurement attributes to get nested inside each observation. 


Datasets are programmatically generated to make sure they are reproducible and scalable. Synthetic data will give the study control over key parameters such as the number of sensors, how frequently data is generated and the duration of the simulation observation period. The idea is to be able to adjust these parameters to produce datasets of different sizes and enable the evaluation of system behavior under varying data volumes. For instance, increasing the number of sensors or sampling frequency creates higher ingestion rates while stretching the time horizon raises the total dataset size.


To make sure that experiments are consistent, the dataset is generated once and then reused across different systems. This removes variability that is introduced by differences in data content and makes sure each system is evaluated under identical conditions. The process of data generation is documented fully including the scripts used and parameters applied to enable reproducibility. 


\section{Workload Design}
This workload design matters to the validity of the evaluation because it defines the operations through which database systems get compared and exercised. While the dataset informs the structure and content of the stored data, the workload specifies the way the data is written, accessed and analyzed. Within the context of time series systems, workload design has to reflect both continuous ingestion of new data and query execution which retrieves and processes historical data. This study’s workload is made to represent two primary classes of operations – query workloads and ingestion workloads. These categories capture the dominant usage patterns in time series applications.


\subsection{Ingestion Workload}
The ingestion workload simulates how sensor-generated data arrives. Data is entered into each database system in sequential batches to reflect the nature of time series data. Instead of entering records individually, the study batches insertion strategies so as to approximate realistic data pipelines where data gets transmitted in grouped intervals instead of as isolated entries. Batch size gets treated like a configurable parameter, letting the evaluation examine how different systems respond to different ingestion patterns.

To make sure that operations are fair across systems, ingestion operations get implemented with the respective database drivers in Python. MongoDB interactions get handled through pymongo, TimescaleDB uses psycopg2 and InfluxDB is accessed via its Python client library. The standardized interaction layer makes sure that all systems get subjected to comparable application-level behavior. Ingestion performance gets measured in terms of throughput, set as the number of records that get successfully inserted per unit time.


\subsection{Query Workload}
Besides ingestion, workload includes queries that are designed to evaluate how the system performs under normal monitoring and analytical scenarios. The queries are made to reflect common operations done on time series data.The query workload is divided into categories based on how they are intended to function. Time-range retrieval is where the system retrieves all observations inside of a specified interval for a given set of sensors or sensor. These queries test how able the system is to scan and filter time-stamped data. 


The second category deals with aggregation queries where statistical summaries like counts, minimums, averages and maximum values get computed over a set time window. These queries are necessary for reporting and monitoring applications and need efficient processing of large data volumes. A third category includes multi-dimensional filtering queries where data gets retrieved based on a combination of metadata attributes and time constraints. These queries typically add a layer of complexity because they require the system to evaluate many filtering conditions at the same time. 

Query workload is executed after there is sufficient volumes of ingested data to make sure that the measurements for performance reflect realistic dataset sizes. The queries are repeated a number of times to account for differences in execution time and the results get averaged to produce stable performance estimates. 



\subsection{\textcolor{red}{Query Execution Approaches}}
\textcolor{red}{This study uses two complementary approaches to provide a comprehensive evaluation of query performance – native query execution and driver-based execution.}


\subsection{Workload Consistency and Control}
For workload design to be effective, it must also factor in consistency across all evaluated systems. For this, the same sequence of operations will be used in each database with equivalent query definitions and identical datasets. Differences in the query syntax between systems get resolved via careful mapping to make sure that each query runs the same logical operation even if the implementation underneath it is different. 

Workload execution is controlled so that there is no interference between query operations and ingestion unless that is explicitly required. Query workloads and ingestion are executed separately to isolate their performance characteristics. This makes sure that the measurements do not get affected by concurrent interactions between workloads. The study is performed this way to make sure the basis for the evaluation is consistent and realistic. 


\section{Experimental Environment and System Configuration}
The environment of the experiment refers to the set of conditions under which all the database systems get deployed and evaluated. When the environment is well defined, it makes performance comparisons reproducible, fair and attributable to characteristics of the system instead of external factors. This section describes the software configuration, hardware setup and deployment approach that is used in the study. 


\subsection{Hardware Environment}
All experiments are run inside a controlled hardware environment to get rid of variability from differences in computational resources. The evaluation is done on a single-node system set with sufficient memory, power and storage capacity to support the defined workloads. This is a deliberate design choice intended to isolate the intrinsic performance characteristics of each database system. Distributed deployments may be common in large-scale production environments but introducing multiple nodes would add a layer of complexity related to synchronization, network communication and cluster management. The focus on a single node configuration makes it possible for the study to ensure that performance differences reflect the efficiency of each system’s storage engine, query processing mechanism and indexing strategy. Hardware specifications including memory capacity, CPU type, storage type and number of cores are documented to make the evaluation reproducible. Where possible, storage devices that have consistent performance characteristics get used to minimize variability in disk I/O behavior. 

\subsection{Software Environment}
Each system is deployed with stable and widely adopted versions to make sure the evaluation is reliable and consistent. The software environment includes the Operating System, database engines and the client libraries used for interaction. InfluxDB, MongoDB and TimescaleDB are installed and configured accordingly. The default configurations become the baseline with minimal adjustments made only according to need to make sure there is the time series functionality or ensure compatibility. That way, the experiment reflects realistic deployment scenarios where systems often get used with standard configurations instead of tuned setups. The experimental framework is run in Python which works like the orchestration layer for query execution, ingestion \textcolor{red}{ and data generation}. Python is chosen because it has an extensive ecosystem of database drivers and it is suitable for automation and scripting. A single programming language across all systems makes sure that there is operational consistency. 


\subsection{Database Configuration}

Even though the default configurations are utilized as the baseline, certain database-specific settings get adjusted to align with the study requirements. For instance, TimescaleDB is configured to enable time-partitioned tables and make sure that data is organized in a way that is consistent with how it will be used. InfluxDB gets configured with necessary measurement definitions and MogoDB uses time series collections to optimize query performance and storage. Indexing strategies get defined consistently across systems to support query workload. In the framework, time-based indexes are also created to enable efficient by timestamp filtering while extra indexes may be applied to sensor identifiers or other attributes that are frequently queried. Care is taken to make sure that no system is unfairly advantaged or disadvantaged and that similar logical structures get used across databases. 

\subsection{Execution Control and Isolation}
To make sure there is consistency, experiments are run under conditions which minimize external process interference. Background applications are limited during testing and the resources of the system monitored to make sure there is no unexpected load that can affect performance measurements. Each experiment is run independently and the database system is reinitialized or reset between runs according to need. This is to prevent caching effects or residual data from influencing the next experiment. Where the experiment requires caching behavior, it is explicitly accounted for in the experimental design. 


\subsection{Considerations for Reproducibility}
Reproducibility is necessary in any empirical evaluation. To support it, all aspects of the experimental environment are recorded including software versions, hardware specifications, execution procedures and configuration settings. The scripts that are used for ingestion, data generation and query execution are preserved and may be reused to replicate the experiments. 





\section{Evaluation Metrics}
This study evaluates database systems using a set of carefully defined performance metrics that capture different aspects of system behavior for time series workloads. Instead of relying on just one measure of performance, the study uses several benchmarks which reflect different operational requirements of large-scale sensor data systems. Every metric is defined the way it operates, with clear procedures for measurement to ensure comparability and consistency across all the studied database technologies. 


\subsection{Ingestion Throughput}
This metric considers the rate at which the system can accept and persist incoming data. It is a measure of the number of records that are successfully inserted into the database per unit time, often expressed in records per second. It is a necessary metric for time series systems where continuous data streams have to be ingested with minimal delay or data loss. To measure it, data is added into each database in controlled batches and the total time used to complete the insertion gets recorded. Throughput gets calculated as the ratio of the number of inserted records to the total time used. The measurements are taken after an initial warm up to make sure that the results reflect a steady state system rather than initialization overhead.


\subsection{Query Latency}
Query latency is a measure of the time it takes for a database system to process a query and return results. It accounts for the system's response time for analytical and monitoring activities, which is a critical factor for applications that depend on the availability of data in a timely manner. Latency is determined by a timing method that measures the beginning and end of a query execution process. For a driver-based execution, the beginning time is measured immediately before the query execution, and the end time is measured after the results are received by the application. The difference between these times is the query latency, which includes the client and network times.


\subsection{Storage Efficiency}
Storage efficiency measures how well a database system uses storage space in handling its time series data. It is defined as the ratio of physical storage space used by the database system for its data to the input data set's size. In calculating storage efficiency, the input data set's storage space is obtained after the ingestion process is complete, including all data structures maintained by the database system, which include indexes and metadata. The input data set's size, on the other hand, is obtained before the ingestion process, allowing for a direct comparison of the input data set and the stored data space. This is a critical factor in a database system, especially in long-term usage, as data is constantly being added over time. A database system that has a higher storage efficiency is ideal for usage in a storage space constraint or a long-term retention scenario.

\subsection{Scalability}
Scalability is a measure used to determine the changes in system performance as the amount of data or intensity of work is increased. Scalability is not defined by a specific numerical value, but by a trend in system performance over different amounts of data. In this particular study, the concept of scalability is analyzed by increasing the size of the data set used in the experiment, repeating the ingestion and query evaluation processes, and analyzing the results obtained. This way, the study can determine the trend in the data ingestion throughput of the system. A system is defined to be scalable if its performance is stable or decreases gracefully in response to increased workload intensity. A sudden decrease in system performance may imply a problem in storage structures, indexing, or resource management.



\subsection{Resource Utilization}
Resource utilization offers an insight into the efficiency of a database system with regards to its ability to effectively utilize resources. The metrics include CPU usage, memory, and disk I/O. Resource utilization is tracked for both ingestion and query execution. High resource usage can imply poor design or configuration of a system, while low resource usage with high performance can imply effective optimization. Resource utilization is not a key aspect of the study but offers useful information for understanding performance results. For instance, two systems can have similar performance but drastically different resource requirements. 




\section{Data Analysis Methods}
The analysis of the results of the experiment needs a structured approach that guarantees reliability, consistency and meaningful interpretation of the observations. This section details the methods used to process, aggregated and compare the data that is collected in the evaluation. The goal is to turn the raw observations into interpretable results while lowering the influence of measurement noise and variability. 


\subsection{Repetition and Averaging}
To make the measurements more reliable, each experiment is performed multiple times in the same environment. This reduces the impact of the system’s transient states and the presence of background processes, among other sources of indeterminism. The results from each run are combined using the average function. This gives a more reliable view of the system’s performance than using a single run’s result. In addition, the range of the results is used to determine the consistency of the system’s performance.



\subsection{Normalization of Results}
To make the results more comparable, some of the results, such as ingestion throughput, are normalized to a standard scale. This is done, for instance, by normalizing the ingestion throughput relative to the best system’s throughput. This is important because it makes it easy to compare the performance of systems whose performance is measured in different ranges.


\subsection{Comparative Analysis}
The main analysis technique is the comparison of performance metrics between various database systems. For instance, for a given performance metric, results are grouped into comparative tables and graphical displays, making it easy to spot performance differences between various database systems. Comparative analysis is done along various axes, including data size, workload types, and complexity. By conducting a multidimensional analysis, it is possible to spot trends in database system behavior, for example, how performance changes with increasing data sizes.



\subsection{Trend Analysis}
Trend analysis is another analysis technique used to analyze performance metrics against changing workload conditions. For example, ingestion throughput can be plotted against data size to understand performance changes with time. Similarly, query latency can be analyzed against different types of queries or time intervals. Understanding trends in performance metrics, makes it possible to understand the performance and reliability of various database systems. For instance, a constant trend in performance metrics is a guarantee of a reliable database system.


\subsection{Handling Variability and Outliers}
It is possible that performance measurements occasionally have outliers resulting from system interruptions, measurement anomalies or resource contention. These outliers get identified through inspection of statistical summaries and repeated measurements. Where possible, the outliers are excluded from the analysis to avoid a distortion of the results. However, the exclusion is carefully justified to make the process transparent. In cases where that variability is significant, it gets reported alongside average values to give a complete picture of system behavior. 


\subsection{Experimental Procedure}
The experimental procedure outlines the steps taken to carry out the evaluation process. The study’s formalization of this process guarantees that all database systems are subject to the same environment, hence a consistent performance measurement process. The process begins with the initialization of the experimental environment. The database systems are set up, initialized, and verified to ensure correct functioning before the experiment. The database systems are reset to a clean state before running the experiment, where necessary, to avoid any effects from previous runs of the experiment.

After initialization, the data set is fed into the database system through a predefined workload. The insertion of data into the database occurs in batches, with the process’s timing being recorded. The process begins with a small warm-up phase to allow the database system to settle before the performance measurement process begins. The initialization process of the database systems is avoided during the experiment to ensure a fair evaluation of the database systems performance.


\section{Threats to Validity and Design Limitations}
The experimental design outlined in this chapter aims to make the experimentation fair and the experimentation environment controlled. However, there are some limitations and potential threats to validity which must be acknowledged. Identifying the limitations is essential to the interpretation of the results accurately and to coming to useful conclusions. 

One potential threat to internal validity comes from any necessary differences in system configuration. Even though efforts are made to make all the configurations standard across database systems, each system is unique and has different tuning parameters and optimization strategies. Going with the default or adjusting configurations minimally may not fully reflect the optimal performance for each system. At the same time, extensive tuning could bring bias into the equation if the tuning is not done consistently. The study mitigates this risk by using a uniform configuration strategy and documenting all used settings. 


Another limitation pertains to the use of synthetic data. Synthetic datasets are preferrable because they allow precise control over the characteristics of the data but they may fail to capture all the complexities of real-world sensor data like missing values, irregular sampling patterns and unpredictable data bursts. Because of this performance observed under the experiment conditions may differ from behavior in real world applications. However, the design of the dataset is set to approximate the key characteristics of time series workloads so that the basis of evaluation is reasonable. 


While there is a rationale for the use of a single-node experimental environment, it introduces extra limitations. The setup allows for controlled comparison of the way the systems behave, but it does not capture the complexities of distributed deployments in which factors like replication, network latency and load balancing play a major role. This may mean that the results won’t fully generalize to large-scale distributed systems. 


Another consideration for the research is measurement accuracy. Timing measurements obtained via application-level instrumentation could include overhead from network communication or client-side processing. Even though this reflects realistic scenarios for use, it may not isolate the internal performance of the database engine. \textcolor{red}{Including native query execution measurements helps to lower this risk by offering an alternative perspective on the system’s performance. }

The selection of database systems may also limit external validity. This study focuses on a subset of the available technologies and the results may not map onto systems that are not included in the evaluation. However, deliberate effort has been made to select databases that represent diverse architectural approaches to allow for meaningful insights into different design paradigms. 

Finally, workload design choices may affect how performance is observed. The workload has been constructed to reflect typical time series operations but variations in ingestion strategies or query patterns could produce different results. This study addresses this by defining the workloads clearly and then applying the consistently across different runs and systems. 







\section{Synthesis of Design Decisions}
The experimental framework design presented in this chapter reflects an intentional effort to balance methodological rigor with practicality. Each component of the design from workload constriction to database selection and performance measurement, is informed by the need to enable fair and meaningful comparison of database systems under time series workloads. 


One central design decision is running the experiments in controlled environments where all systems get evaluated under identical conditions. This is done to make sure that performance differences are directly attributable to system behavior instead of external variability. Using a single-node configuration helps to anchor this objective by isolating the characteristics intrinsic to each database system even though it adds limitations to distributed scalability analysis. 


How the database systems are selected reflects an emphasis on architectural diversity. The study includes relational, document-oriented approaches and native time series to examine the way different design philosophies influence performance. This diversity improves the analytical value of the study. At the same time, the workload and dataset designs are structured to reflect realistic usage scenarios while guaranteeing control over experimental variables. The use of synthetic data provides flexibility and reproducibility and defining query and ingestion workloads ensures that each system is evaluated under representative conditions. Separating query and ingestion phases makes clear analysis possible for each performance dimension without any interferences. 


Performance metrics are defined in ways that capture many aspects of system behavior including scalability, throughput, latency, storage efficiency and resource utilization. This approach acknowledges that no one metric is enough to fully characterize system performance and allows for the comparison to be more nuanced.


The experimental procedure and the methods for data analysis are meant to ensure consistency, transparency and repeatability. Repeated measurements, systematic analysis techniques and structured data collection contribute to reliable results and support the validity of the conclusions this study draws. Taken together, these decisions for design create a coherent and robust framework to evaluate database systems for time series data management. The framework also provides a basis for the experimental results provided in subsequent chapters. 

